\scanpage{77}
\mspage{- 17 -}

$|z_i|\leq r_i$ ; soit $M$ la borne supérieure des normes des termes de la série pour $|z_i|=r_i$ $(1\leq i\leq p)$ ; on a donc

\[
\left\|\handwritten{\unsure{c|\mathbb C|C}}_{n_1n_2..n_p}\right\|\leq M/r_1^{n_1}..r_p^{n_p}
\]

ce qui montre que la série entière $(M\,(\frac{x_1}{r_1})^{\handwritten{n}_1}\ldots(\frac{x_p}{r_p})^{\handwritten{n}_p})$ est une majorante de la série donnée ; on observera que la somme de cette majorante est $M/(1-\frac{x_1}{r_1})(1-\frac{x_2}{r_2})\ldots(1-\frac{x_p}{r_p})$ pour $|x_i|<r_i$ $(1\leq i\leq p)$. Une autre majorante est la série obtenue par substitution de $\handwritten{\sum_{i=1}^p}\frac{x_i}{r_i}$ à $u$ dans la série géométrique $(Mu^n)$ ; en effet, il est immédiat que le coefficient de $x_1^{n_1}..x_p^{n_p}$ dans cette série est le produit d'un entier $\geq1$ et de $M/r_1^{n_1}..r_p^{n_p}$ ; cette nouvelle majorante a pour somme $M/(1-\handwritten{\sum_{i=1}^p}\frac{x_i}{r_i})$ dans un voisinage assez petit de l'origine.

\uline{2}. \uline{Résolution des équations analytiques}.

THEOREME 1 .- \uline{Soient} $f_i(x_1,..,x_p\,;\,y_1,y_2,\ldots,y_p)$ $(1\leq i\leq q)$ $q$ \uline{fonctions holomorphes scalaires définies dans une partie ouverte} $A$ \uline{de} $K^{p+q}$. \uline{Soit} $(a_1,\ldots,a_p\,;\,b_1,..,b_q)$ \uline{un point de} $A$ \uline{tel que} $f_i(a_1,..a_p\,;\,b_1,..,b_q)=0$ \uline{pour} $1\leq i\leq q$, \uline{et que le déterminant fonctionnel} $\handwritten{\Delta}=\det(\handwritten{\partial} f_i/\handwritten{\partial} y_j)$ \uline{soit} $\handwritten{\ne}0$ \uline{en ce point}. \uline{Alors il existe un voisinage} $V$ \uline{de} $(a_1,..,a_p)$ \uline{dans} $K^p$ \uline{et un système de} $q$ \uline{fonctions scalaires} $u_i(x_1,..,x_p)$ $(1\leq i\leq q)$ \uline{holomorphes dans} $V$, \uline{satisfaisant aux conditions} $u_i(a_1,..,a_p)=b_i$ $(1\leq i\leq q)$ \uline{et}

\[
f_i(x_1,\ldots,x_p\,;\,u_1(x_1,..,x_p),..,u_q(x_1,..,x_p))=0\qquad(1\leq i\leq q)\tag*{(1)}
\]

\uline{en tout point de} $V$. \uline{En outre}, \uline{si} $v_i$ $(1\leq i\leq q)$ \uline{est un système de} $q$ \uline{fonctions continues au point} $(a_1,\ldots,a_p)$, \uline{définies dans un voisinage de ce point}, \uline{et satisfaisant aux conditions} $v_i(a_1,..,a_p)=b_i$ $(1\leq i\leq q)$ \uline{et}

% handwritten notes in the opposite direction to the typescript
\handwritten{$\deleted{\unsure{\mathbb Z_2|2}}\ \mathbb Z_p($}

\handwritten{$Q(i)^{\deleted{\unsure{*}}}\qquad T$}

\handwritten{$\unsure{B^o_C|B^0_C|B^o_c|\mathcal{B}_c}(\mathbb Z)$}

\handwritten{$-1$ \inserted{pas un} carré mod $p$\qquad $p\equiv-1\qquad(4)$}

\handwritten{$p\qquad\left[\genfrac{}{}{0pt}{}{-1}{p}\right]=-1\qquad p$}

\handwritten{\unsure{TORS|TORs}(}
% archive sheet number: 76

